\section{Related Work}
\noindent\textbf{Memory-Augmented LLM Agents.}
External memory has become a core component of LLM agents that operate over long horizons.
Prior work improves long-term memory from several directions, including memory architecture~\cite{packer2023memgpt,zhong2024memorybank}, memory organization and consolidation~\cite{xu2025mem,fang2025lightmem}, and memory retrieval~\cite{du2025memr}.
These methods substantially improve individual stages of the memory pipeline, but they primarily optimize storage, organization, or retrieval in isolation.
Our work is inherently different from existing work: \method jointly coordinates memory construction and iterative retrieval, and converts utilization failures into direct repair signals for the memory bank.
% By contrast, \method addresses a broader scope: it coordinates both construction and retrieval along the forward path of the memory cycle, and further converts utilization failures into direct repair signals for the memory bank along the backward path. 
Full version is in Appendix~\ref{appendix:full_related_works}.

% \noindent\textbf{Self-Evolution for LLM Agents.}
% Recent work improves agent behavior through output-level 
% refinement~\cite{madaan2023self}, 
% experience accumulation~\cite{zhao2024expel,wang2023voyager}, 
% and policy optimization for memory~\cite{shen2026membuilder,yan2025memory}.
% These approaches improve responses, maintain separate experience 
% stores, or optimize memory policies through training, but do not 
% directly repair the memory bank during construction.
% \method fills this gap by verifying the current memory with probe 
% QA pairs and converting failures into repair actions before the 
% memory is committed. Full version is in Appendix~\ref{appendix:self_evolve_LLM_related_works}.